\documentclass{article}
\usepackage{amsmath, amssymb, amsthm}
\usepackage{hyperref}
\usepackage{geometry}
\geometry{margin=1in}

\title{Erd\H{o}s Problem \#940}
\date{Last updated: 2025-08-31}
\author{Source: erdosproblems.com}

\begin{document}
\maketitle

\noindent\textbf{Status:} OPEN \quad \textbf{No prize}

\noindent\textbf{Tags:} number theory

\medskip

\noindent\textbf{Problem Statement:}

Let $r\geq 3$. A number $n$ is $r$-powerful if for every prime $p$ which divides $n$ we have $p^r\mid n$. Are there infinitely many integers which are not the sum of at most $r$ many $r$-powerful numbers? Does the set of integers which are the sum of at most $r$ $r$-powerful numbers have density $0$?




Erd\H{o}s \cite{Er76d} claims that the claim that the set has density $0$ is 'easy' for $r=2$ (a potential 'easy argument' is given in the comments by Tao; this was first proved in the literature by Baker and Br\"{u}dern \cite{BaBr94}). For $r=3$ it is not even known if those integers which are the sum of at most three cubes has density $0$. In the Oberwolfach problem book this is recorded in 1986 as a problem of Erd\H{o}s and Ivi\'{c}. In \cite{Er76d} Erd\H{o}s claims that 'a simple counting argument' implies that there are infinitely many integers which are not the sum of at most $r$ many $r$-powerful numbers, but Schinzel pointed out he made a mistake. Heath-Brown \cite{He88} has proved that all large numbers are the sum of at most three $2$-powerful numbers, see [941] . See also [1081] for a more refined question concerning the case $r=2$, and [1107] for the case of $r+1$ summands.




References

[BaBr94] Baker, R. C. and Br\"udern, J., On sums of two squarefull numbers . Math. Proc. Cambridge Philos. Soc. (1994), 1--5.

[Er76d] Erd\H{o}s, P., Problems and results on number theoretic properties of consecutive integers and related questions . Proceedings of the Fifth Manitoba Conference on Numerical Mathematics (Univ. Manitoba, Winnipeg, Man., 1975) (1976), 25-44.

[He88] Heath-Brown, D. R., Ternary quadratic forms and sums of three square-full numbers . (1988), 137--163.

\noindent\textbf{Related OEIS sequences:} possible


\bigskip
\noindent\small{Source: \url{https://www.erdosproblems.com/940}}
\end{document}
